\documentclass{article}
\usepackage{amsmath, amssymb, physics}

\title{Synthetic Paper for Benchmark B003-003}
\author{A. Example \\ Example University}
\date{\today}

\begin{document}
\maketitle

\section{Introduction}
This is a synthetic paper for benchmarking the math paper replay auditor.

\section{Main Results}
\begin{theorem}[Main Theorem]
\label{thm:main}
Let $f(x) = e^x$. Then for all $x \in \mathbb{R}$, we have
\begin{equation}
\label{eq:main}
f'(x) = f(x).
\end{equation}
\end{theorem}

\begin{proof}
By definition of the derivative,
\begin{equation}
\label{eq:deriv}
f'(x) = \lim_{h \to 0} \frac{f(x+h) - f(x)}{h} = \lim_{h \to 0} \frac{e^{x+h} - e^x}{h} = e^x \lim_{h \to 0} \frac{e^h - 1}{h} = e^x.
\end{equation}
\end{proof}

\begin{lemma}[Auxiliary Lemma]
\label{lem:aux}
For all $x \in \mathbb{R}$, $e^x > 0$.
\end{lemma}

\begin{corollary}
\label{cor:result}
From Theorem~\ref{thm:main} and Lemma~\ref{lem:aux}, $f'(x) > 0$ for all $x \in \mathbb{R}$.
\end{corollary}

\section{Numerical Examples}
\begin{table}[h]
\centering
\begin{tabular}{|c|c|}
\hline
$x$ & $f(x)$ \\
\hline
0 & 1.00000 \\
1 & 2.71829 \\
2 & 7.38906 \\
\hline
\end{tabular}
\caption{Numerical example for $f(x) = e^x$}
\end{table}

\end{document}
